\documentclass[11pt]{article}
% Shared preamble for all ML-RESEARCH papers.
% Papers load it with: % Shared preamble for all ML-RESEARCH papers.
% Papers load it with: % Shared preamble for all ML-RESEARCH papers.
% Papers load it with: \input{../../../papers-common/preamble.tex}
\usepackage[utf8]{inputenc}
\usepackage[T1]{fontenc}
\usepackage{lmodern}
\usepackage{amsmath,amssymb,amsthm}
\usepackage{graphicx}
\usepackage{booktabs}
\usepackage{hyperref}
\usepackage[margin=1.1in]{geometry}

\newtheorem{theorem}{Theorem}
\newtheorem{proposition}{Proposition}
\newtheorem{definition}{Definition}
\newtheorem{remark}{Remark}

\newcommand{\R}{\mathbb{R}}
\newcommand{\E}{\mathbb{E}}
\newcommand{\Prob}{\mathbb{P}}
\newcommand{\ip}[2]{\langle #1, #2 \rangle}

\usepackage[utf8]{inputenc}
\usepackage[T1]{fontenc}
\usepackage{lmodern}
\usepackage{amsmath,amssymb,amsthm}
\usepackage{graphicx}
\usepackage{booktabs}
\usepackage{hyperref}
\usepackage[margin=1.1in]{geometry}

\newtheorem{theorem}{Theorem}
\newtheorem{proposition}{Proposition}
\newtheorem{definition}{Definition}
\newtheorem{remark}{Remark}

\newcommand{\R}{\mathbb{R}}
\newcommand{\E}{\mathbb{E}}
\newcommand{\Prob}{\mathbb{P}}
\newcommand{\ip}[2]{\langle #1, #2 \rangle}

\usepackage[utf8]{inputenc}
\usepackage[T1]{fontenc}
\usepackage{lmodern}
\usepackage{amsmath,amssymb,amsthm}
\usepackage{graphicx}
\usepackage{booktabs}
\usepackage{hyperref}
\usepackage[margin=1.1in]{geometry}

\newtheorem{theorem}{Theorem}
\newtheorem{proposition}{Proposition}
\newtheorem{definition}{Definition}
\newtheorem{remark}{Remark}

\newcommand{\R}{\mathbb{R}}
\newcommand{\E}{\mathbb{E}}
\newcommand{\Prob}{\mathbb{P}}
\newcommand{\ip}[2]{\langle #1, #2 \rangle}


\title{Loss-Based Training of the Linear Discriminator:\\ Experimental Writeup}
\author{Yuval Lavie}
\date{\today}

\begin{document}
\maketitle

\begin{abstract}
We train the pure-Python linear discriminator of paper~01 on synthetic
two-Gaussian data by full-batch gradient descent on the logistic loss. All
numbers and figures in this document are generated from the experiment
tracker (\texttt{tracking.db}) by \texttt{generate\_assets.py}; nothing is
typed in by hand. To reproduce: run the configs under
\texttt{experiments/}, then \texttt{mlr paper} on this directory.
\end{abstract}

\section{Setup}

The dataset is two isotropic Gaussians in $\R^2$ with unit variance, centered
$\pm\,\mathrm{sep}/2$ apart on the first coordinate, $n = 1000$ points,
balanced classes, $20\%$ held out for testing. The model is trained with
full-batch gradient descent (learning rate $0.5$, $300$ epochs) on the
logistic loss, as derived in paper~01.

\section{Results}

\begin{table}[h]
  \centering
  \input{../../results/tables/accuracy.tex}
  \caption{Test accuracy per experiment, pulled from the latest finished run
  of each config in the tracker.}
  \label{tab:accuracy}
\end{table}

Wider class separation makes the problem easier, and test accuracy rises
accordingly (Table~\ref{tab:accuracy}); the residual error at small
separation is the irreducible overlap of the two Gaussians, which no
hypothesis in the linear class can remove.

\begin{figure}[h]
  \centering
  \includegraphics[width=0.75\textwidth]{../../results/figures/loss-curve.pdf}
  \caption{Training loss per epoch for each experiment. The easier
  (wider-separation) problem drives the logistic loss lower, consistent with
  the likelihood interpretation of paper~01.}
  \label{fig:loss}
\end{figure}

\end{document}
